(1) 令圆柱径向坐标为 $r=\sqrt{x^2+y^2}$。在导体外的等效解是实电荷 $q$ 放在 $(0,0,h)$、镜像电荷 $-q$ 放在 $(0,0,-h)$，故 $\Phi(r,z)=\dfrac{q}{4\pi\varepsilon_0}\left[\dfrac1{\sqrt{r^2+(z-h)^2}}-\dfrac1{\sqrt{r^2+(z+h)^2}}\right]$，适用域为 $z>0$。在 $z=0$ 两项相消，满足接地边界，唯一性定理保证此解正确。
(2) 导体向真空的法向为 $+\hat z$，电场法向 $E_z(0^+)=-\partial_z\Phi|_{0^+}=-qh/[2\pi\varepsilon_0(r^2+h^2)^{3/2}]$，所以 $\sigma(r)=\varepsilon_0E_z=-qh/[2\pi(r^2+h^2)^{3/2}]$。面积积分 $Q_{\rm ind}=\int_0^\infty\sigma(r)2\pi r\,dr=-qh\int_0^\infty r(r^2+h^2)^{-3/2}dr=-q$。
(3) 导体对真实电荷的作用等于镜像电荷在真实电荷位置产生的场力。两电荷的镜像距离是 $2h$，因此 $\boldsymbol F=-\hat z\,q^2/[4\pi\varepsilon_0(2h)^2]=-\hat z\,q^2/(16\pi\varepsilon_0h^2)$。这是指向导体的吸引力，不能在力的公式中再额外乘二。
(4) 由 $F_h=-dU/dh=-q^2/(16\pi\varepsilon_0h^2)$，且取 $U(\infty)=0$，积分得 $U(h)=-q^2/(16\pi\varepsilon_0h)$。若把镜像当成第二颗真实可独立移动电荷，直接得到的成对相互作用能 $-q^2/(8\pi\varepsilon_0h)$ 会大一倍；镜像随真实电荷位置而改变，正确的固定接地边界充电能只取镜像对相互作用能的一半。
